\textbf{1. Kiến thức cần nhớ}
\begin{itemize}
\item \textbf{Tần số} của một giá trị là số lần xuất hiện của giá trị đó trong mẫu dữ liệu, kí hiệu là $m$ (hoặc $f$).
\item Tổng các tần số của tất cả các giá trị bằng cỡ mẫu (tổng số phần tử dữ liệu $N$):
$$\boxed{m_1 + m_2 + \dots + m_k = N}$$
\item \textbf{Bảng tần số}: Gồm hai dòng (hoặc hai cột): dòng trên ghi các giá trị khác nhau $x_i$, dòng dưới ghi các tần số tương ứng $m_i$.
\item \textbf{Biểu đồ tần số}: Biểu đồ cột hoặc biểu đồ đoạn thẳng biểu diễn tương ứng giữa giá trị và tần số.
\end{itemize}

\textbf{2. Phương pháp giải các dạng}

%% DANG 9D7?22-1
\textbf{Dạng 1. Lập bảng tần số}
\begin{enumerate}
\item Liệt kê tất cả các giá trị khác nhau xuất hiện trong mẫu số liệu theo thứ tự từ nhỏ đến lớn.
\item Đếm số lần xuất hiện (tần số $m_i$) của từng giá trị.
\item Kiểm tra lại: tổng tất cả các tần số phải bằng đúng cỡ mẫu $N$.
\end{enumerate}

%% DANG 9D7?22-2
\textbf{Dạng 2. Vẽ và đọc biểu đồ tần số}
\begin{enumerate}
\item Trục ngang biểu diễn các giá trị $x_i$, trục đứng biểu diễn tần số $m_i$.
\item Vẽ cột (hình chữ nhật) hoặc đoạn thẳng nối các đỉnh tương ứng $(x_i; m_i)$.
\item Đọc biểu đồ: Gióng từ cột/điểm sang trục đứng để xác định tần số của từng giá trị; tìm giá trị có tần số lớn nhất (mốt).
\end{enumerate}
